(1) 令 \(U(t)=e^{-i\Omega t\sigma_z/2}\) 并写 \(|\psi(t)\rangle=U(t)|\chi(t)\rangle\)。把旋转项移到框架内，\(H_{\rm rot}=U^\dagger HU-i\hbar U^\dagger\dot U\)。利用 \(U\sigma_xU^\dagger=\cos\Omega t\,\sigma_x+\sin\Omega t\,\sigma_y\)，及 \(-i\hbar U^\dagger\dot U=-\hbar\Omega\sigma_z/2\)，得到 \(H_{\rm rot}=-(\hbar/2)[\omega_0\sin\theta\,\sigma_x+(\omega_0\cos\theta+\Omega)\sigma_z]\)。与原哈密顿量相比新增项的符号由旋转变换和时间惯例决定，不得随意丢失。

(2) 定义有效频率 \(\omega_{\rm eff}=\sqrt{\omega_0^2+\Omega^2+2\omega_0\Omega\cos\theta}\)，旋转框架的两本征能为 \(\mp\hbar\omega_{\rm eff}/2\)，间隔为 \(\hbar\omega_{\rm eff}\)。初态 Bloch 方向 \(\mathbf n(0)=(\sin\theta,0,\cos\theta)\)，有效场方向 \(\mathbf h=(\omega_0\sin\theta,0,\omega_0\cos\theta+\Omega)/\omega_{\rm eff}\)。二者叉积模长 \(|\mathbf n\times\mathbf h|=\Omega\sin\theta/\omega_{\rm eff}\)，故夹角 \(\beta\) 满足 \(\sin\beta=\Omega\sin\theta/\omega_{\rm eff}\)，这就是非绝热混合的几何幅度。

(3) 时间无关旋转框架演化算符为 \(\exp[-iH_{\rm rot}t/\hbar]=\cos(\omega_{\rm eff}t/2)I+i(\mathbf h\cdot\boldsymbol\sigma)\sin(\omega_{\rm eff}t/2)\)。相对于初态正交自旋方向的跃迁概率为 \(P_{\rm flip}(t)=\sin^2\beta\,\sin^2(\omega_{\rm eff}t/2)\)，所以一周期后的精确结果 \(P_{\rm flip}(T)=[\Omega^2\sin^2\theta/\omega_{\rm eff}^2]\sin^2(\pi\omega_{\rm eff}/\Omega)\)。由于旋转算符 \(U(T)=-I\) 只引入共同的自旋号因子，不改变这一概率。若 \(\Omega/\omega_0\to0\)，非绝热跃迁包络至多是 \(O[(\Omega\sin\theta/\omega_0)^2]\)，因此绝热条件可取 \(\Omega\sin\theta\ll\omega_0\)。

(4) 缓慢绕转时磁场单位矢量扫出的立体角为 \(\mathcal A=2\pi(1-\cos\theta)\)。沿瞬时 \(\boldsymbol\sigma\cdot\mathbf n\) 正向的自旋本征态积累 Berry 相位 \(\gamma_+=-\mathcal A/2=-\pi(1-\cos\theta)\)，反向态为 \(\gamma_-=\mathcal A/2\)，均按模 \(2\pi\) 理解。它与动力相位 \(-\hbar^{-1}\int E_\pm dt\) 源自不同几何机制；单一态的整体相位不可单独观测，但两路径或两自旋支之间的相位差可通过干涉测量。有限转速下还应使用前述精确结果而非仅套用 Berry 相位。
